\chapter{Build and run guide}\label{app:build}

This appendix collects the commands to build the port and run its tests, in the order a new contributor needs them.
The authoritative versions are in the repository: \code{port/agent/ENV_H100.md} (machine setup),
\code{port/agent/WORKFLOW.md} (the work loop) and the usage lines at the top of each script.

\section{Without a GPU}

The first three rungs of the routine ladder (\cref{ch:port}) need only gfortran and Python:
\begin{lstlisting}[language={}]
bash port/gates/static.sh                         # guards: arith_guard, kernel_lint, scope, ...
bash port/h100/compile_one.sh gnu-gpu <file>      # one file, GPU view (gfortran -fopenacc)
bash port/h100/compile_one.sh gnu-ref <file>      # one file, CPU view
bash port/h100/harness.sh <file> <routine> --builds gnu-ref,gnu-gpu
CONTAINER=none NETCDF=/opt/netcdf bash port/gates/cpu_verify.sh   # whole-model CPU check
\end{lstlisting}
\code{cpu_verify.sh} builds both gfortran views and the CPU view of the base commit from clean directories, runs the
smoke case S-3M with each, and requires the GPU view to equal the CPU view bit for bit. It also runs the reference unit
tests (templates, mutants, the limiter, the ozone interpolation, KISS).

\section{Setting up the GPU machine}

Everything runs in NVIDIA's HPC SDK container, without root:
\begin{lstlisting}[language={}]
cp port/h100/env.sh port/h100/env.local.sh        # set WORK, CASE_INPUTS, CONTAINER, IMAGE, CPU_RANKS, GPU_ID
bash port/h100/setup_toolchain.sh image           # pull the pinned NVHPC image
bash port/h100/setup_toolchain.sh deps            # tcsh, m4, HDF5 and netCDF built with nvfortran
bash port/h100/setup_toolchain.sh python          # numpy, netCDF4, mpmath for the host tools
bash port/h100/setup_toolchain.sh check           # must end with "toolchain: PASS"
\end{lstlisting}
The compiler version is pinned for the whole port: the arithmetic tests and probes are valid only for the version
they ran with.

\section{Builds}

\begin{lstlisting}[language={}]
bash port/h100/build.sh <mode> --commit REV       # clean build of a commit, reused if it exists
bash port/h100/build.sh <mode> --worktree         # incremental build of the working tree
\end{lstlisting}
\begin{center}
\small
\begin{tabular}{lp{0.7\textwidth}}
\toprule
Mode & What \\
\midrule
\code{cpu-ref} & CPU-REF: nvfortran, reproducible flags, \code{REPRO_MATH}, \code{WRF_POOL} (\cref{ch:iobuild}) \\
\code{gpu-repro} & GPU-REPRO: as CPU-REF, plus \code{-acc=gpu -cuda -gpu=cc80,cc90,nofma,noflushz} and \code{WRF_GPU} \\
\code{gpu-debug} & GPU-REPRO with debugging information and line numbers \\
\code{gnu-ref}, \code{gnu-gpu} & the two gfortran views, for checks without NVHPC or a GPU \\
\bottomrule
\end{tabular}
\end{center}
Options: \code{--fine} adds the fine trace checkpoints; \code{--uninit nan|zero} (gfortran) builds the two halves of
T-UNINIT. A clean build of WRF takes 30 to 90 minutes, an incremental build after editing one file a few minutes.

\section{The dev case and its references}

\begin{lstlisting}[language={}]
bash port/h100/dev_case.sh make                   # checks the inputs' md5s, builds eaton_small
nohup bash port/h100/dev_case.sh reference > $WORK/dev_reference.log 2>&1 &   # CPU-REF references, hours
bash port/h100/dev_case.sh status
\end{lstlisting}

\section{Windows and comparisons}

\begin{lstlisting}[language={}]
bash port/h100/window.sh <build dir> W-20 [VAR=value ...]   # one window; prints the run directory
bash port/h100/compare.sh <run A> <run B>          # traces and files, bit for bit
bash port/gates/t_ab.sh <route> W-20               # route on against route off
bash port/gates/t_trace.sh W-20                    # GPU-REPRO against CPU-REF
bash port/gates/t_reg20.sh                         # the per-commit regression
\end{lstlisting}
The windows are defined in \code{port/h100/windows.txt} (\cref{app:tests}). Runs are reused while the binary is the same.

\section{Environment variables}

\begin{center}
\small
\begin{tabular}{lp{0.65\textwidth}}
\toprule
Variable & Effect \\
\midrule
\code{WRF_GPU_OFF=r1,r2} & run these routes on the host (T-AB) \\
\code{WRF_GPU_ONLY=r1,r2} & run only these routes on the device \\
\code{WRF_GPU_CALLCHECK=r:n} & check the first $n$ calls of route $r$, device against host (\cref{alg:callcheck}) \\
\code{WRF_BITTRACE=0|1|2} & tracer off; end of every step; every checkpoint \\
\code{WRF_BITTRACE_FROM}, \code{_TO} & range of steps to trace \\
\code{WRF_GPU_TIMING=1} & time per NVTX range, per simulated hour \\
\code{WRF_RRTMG_BATCH=n} & columns per RRTMG batch (default 4096) \\
\code{NV_ACC_CUDA_STACKSIZE} & device stack per thread, for column physics \\
\code{ACC_PROFLIB} & load an OpenACC profiling library (\cref{ch:prof}) \\
\bottomrule
\end{tabular}
\end{center}

\section{The work loop}

\begin{lstlisting}[language={}]
python3 port/tools/workbook.py resume              # where the work stands
python3 port/tools/workbook.py next                # what to do
python3 port/tools/ref.py <kernel id>              # the CPU lines to port
# edit under #ifdef WRF_GPU, following port/agent/CODING_STANDARD.md
bash port/gates/static.sh
bash port/h100/compile_one.sh gpu-repro <file> --minfo
bash port/h100/harness.sh <file> <routine>
bash port/gates/t_ab.sh <route> W-20
bash port/gates/t_trace.sh W-20
python3 port/tools/workbook.py set <kernel> done --commit <sha> --tests "..."
\end{lstlisting}

\section{Building this book}

\begin{lstlisting}[language={}]
cd book && make                                     # latexmk -pdf; needs TeX Live
python3 book/tools/kernels_tex.py                   # regenerate Appendix A from kernels.csv
\end{lstlisting}
On every push that changes \code{book/}, GitHub Actions builds the PDF and keeps it as the artifact
\code{wrf-gpu-book}; the build fails on any undefined reference or citation.
